\section{Causal Dynamical Triangulations -- 2D}

Causal Dynamical Triangulations (CDT) is a non-perturbative approach to quantum gravity that constructs spacetime from fundamental simplicial building blocks while preserving a causal (Lorentzian) structure.

\subsection{The Path Integral}

The gravitational path integral is regularized by summing over causal triangulations $T$:

\begin{equation}
Z = \sum_{T} \frac{1}{C_T} e^{-S_{\text{Regge}}[T]}
\end{equation}

\begin{itemize}
\item $C_T$ -- symmetry factor of triangulation $T$
\item $S_{\text{Regge}}$ -- Regge action (discrete Einstein--Hilbert)
\end{itemize}

\subsection{Regge Action in 2D}

In 1+1 dimensions the Regge action simplifies to:

\begin{equation}
S = \lambda N_2 - \mu N_0
\end{equation}

\begin{itemize}
\item $\lambda$ -- cosmological constant coupling
\item $N_2$ -- number of triangles
\item $N_0$ -- number of vertices
\end{itemize}

\subsection{Monte Carlo Moves}

The triangulation is updated via Pachner moves that preserve the causal structure:

\begin{itemize}
\item $(2,2)$ flip -- exchanges the diagonal of two adjacent triangles
\item $(2,4)$ insertion -- splits a triangle by inserting a vertex
\item $(4,2)$ removal -- removes a vertex and merges triangles
\end{itemize}

\subsection{Observables}

\param{$d_H$}{Hausdorff dimension -- measures how volume scales with geodesic distance}
\param{$d_s$}{Spectral dimension -- extracted from the return probability of a diffusion process}
\param{$n(t)$}{Volume profile -- number of triangles per time slice}
